% ==========================
% 标题与目录与排版格式（最终清洁版本）
% ==========================

% 目录样式：章“第X章”，节“X、”，小节“1.”，段落“(1)”
\ExplSyntaxOn
\csname tl_new:N\endcsname \l__toc_label_tl
\seq_new:N \l__toc_label_seq
\cs_new_protected:Npn \__toc_prepare_label: {
  \csname tl_set:Nx\endcsname \l__toc_label_tl { \thecontentslabel }
  \csname tl_if_blank:VTF\endcsname \l__toc_label_tl
    { \seq_clear:N \l__toc_label_seq }
    {
      \seq_set_split:NnV \l__toc_label_seq {.} \l__toc_label_tl
      \iow_term:x { label~=~\tl_use:N \l__toc_label_tl }
      \seq_map_inline:Nn \l__toc_label_seq { \iow_term:x { ~component~=~##1 } }
    }
}
\cs_new:Npn \__toc_component_or_label:n #1 {
  \int_compare:nNnTF {#1} > {\seq_count:N \l__toc_label_seq}
    {
      \iow_term:n { fallback~label }
      \l__toc_label_tl
    }
    {
      \seq_item:Nn \l__toc_label_seq { #1 }
    }
}
\cs_new_protected:Npn \__toc_if_label:nnn #1#2#3 {
  \csname tl_if_blank:VTF\endcsname \l__toc_label_tl { #3 } { #1#2 }
}
\cs_new_protected:Npn \__toc_section_entry:n #1 {
  	ypeout{section label raw = \meaning\thecontentslabel}
  \__toc_prepare_label:
  \__toc_if_label:nnn
    { 第\zhnumber{\__toc_component_or_label:n { 1 }}章\quad }
    { #1 }
    { #1 }
}
\cs_new_protected:Npn \__toc_subsection_entry:n #1 {
  	ypeout{subsection label raw = \meaning\thecontentslabel}
  \__toc_prepare_label:
  \__toc_if_label:nnn
    { \zhnumber{\__toc_component_or_label:n { 2 }}、\quad }
    { #1 }
    { #1 }
}
\cs_new_protected:Npn \__toc_subsubsection_entry:n #1 {
  	ypeout{subsubsection label raw = \meaning\thecontentslabel}
  \__toc_prepare_label:
  \__toc_if_label:nnn
    { \__toc_component_or_label:n { 3 }.\quad }
    { #1 }
    { #1 }
}
\cs_new_protected:Npn \__toc_paragraph_entry:n #1 {
  	ypeout{paragraph label raw = \meaning\thecontentslabel}
  \__toc_prepare_label:
  \__toc_if_label:nnn
    { (\__toc_component_or_label:n { 4 })\quad }
    { #1 }
    { #1 }
}
\ExplSyntaxOff
\DeclareRobustCommand{\tocSectionEntry}[1]{\__toc_section_entry:n {#1}}
\DeclareRobustCommand{\tocSectionEntryNoLabel}[1]{#1}
\DeclareRobustCommand{\tocSubsectionEntry}[1]{\__toc_subsection_entry:n {#1}}
\DeclareRobustCommand{\tocSubsectionEntryNoLabel}[1]{#1}
\DeclareRobustCommand{\tocSubsubsectionEntry}[1]{\__toc_subsubsection_entry:n {#1}}
\DeclareRobustCommand{\tocSubsubsectionEntryNoLabel}[1]{#1}
\DeclareRobustCommand{\tocParagraphEntry}[1]{\__toc_paragraph_entry:n {#1}}
\DeclareRobustCommand{\tocParagraphEntryNoLabel}[1]{#1}

\csname titlecontents\endcsname{section}[0em]{\normalfont\normalsize}{\tocSectionEntry}{\tocSectionEntryNoLabel}{\titlerule*[0.5pc]{.}\contentspage}
\csname titlecontents\endcsname{subsection}[2em]{\normalfont\normalsize}{\tocSubsectionEntry}{\tocSubsectionEntryNoLabel}{\titlerule*[0.5pc]{.}\contentspage}
\csname titlecontents\endcsname{subsubsection}[3.5em]{\normalfont\normalsize}{\tocSubsubsectionEntry}{\tocSubsubsectionEntryNoLabel}{\titlerule*[0.5pc]{.}\contentspage}
\csname titlecontents\endcsname{paragraph}[5em]{\normalfont\normalsize}{\tocParagraphEntry}{\tocParagraphEntryNoLabel}{\titlerule*[0.5pc]{.}\contentspage}

% 标题格式
\titleformat{\section}{\centering\normalfont\zihao{-3}}{第\chinese{section}章}{1em}{}
\titlespacing{\section}{0pt}{1\baselineskip}{1\baselineskip}
\titleformat{\subsection}{\normalfont\zihao{4}}{\chinese{subsection}、}{1em}{}
\titlespacing{\subsection}{0pt}{0.8\baselineskip}{0.6\baselineskip}
\titleformat{\subsubsection}{\normalfont\zihao{-4}}{\arabic{subsubsection}.}{0.8em}{}
\titlespacing{\subsubsection}{0pt}{0.7\baselineskip}{0.4\baselineskip}
\titleformat{\paragraph}{\normalfont\zihao{-4}}{(\arabic{paragraph})}{0.8em}{}
\titlespacing{\paragraph}{0pt}{0.6\baselineskip}{0.2\baselineskip}

% 进入新的 subsubsection 时重置 paragraph 计数
\makeatletter
\@addtoreset{paragraph}{subsubsection}
\makeatother

% 编号与目录深度
\setcounter{secnumdepth}{4}
\setcounter{tocdepth}{3}

% 表格与段落基础格式
\renewcommand{\arraystretch}{1.5}
\setlength{\parindent}{2em}

% 目录标题（居中：目 录）
\renewcommand{\contentsname}{ \centering{目\hspace{1em}录}}

% 自定义辅助命令
\ExplSyntaxOn
\NewDocumentCommand {\hfillspace}{ > { \SplitList { ~ } } m }{ \tl_map_inline:nn {#1} { ##1\hfill } }
\ExplSyntaxOff
\newcounter{myline}
\AtBeginEnvironment{tabular}{\setcounter{myline}{0}}

% 下划线固定宽度命令
\newcommand{\fixedunderline}[2][6cm]{\underline{\makebox[#1][c]{#2}}}

% 代码高亮设置
\lstset{
  language=Matlab,
  basicstyle=\ttfamily,
  numbers=left,
  numberstyle=\tiny,
  stepnumber=1,
  numbersep=5pt,
  showspaces=false,
  showstringspaces=false,
  showtabs=false,
  frame=single,
  tabsize=2,
  breaklines=true,
  breakatwhitespace=true,
  captionpos=b,
  keepspaces=true,
  columns=flexible,
}

